% GENERATED FILE. Edit canonical lessons, then run npm run generate:books.
% canonical-source-hash: fnv1a64:8b1e33f1b693b619
% canonical-lessons: LA-C494-odi, LA-C494-odiosus, LA-C494-fastidio, LA-C494-fastidium, LA-C494-taedet

\chapter{Hatred and disgust}
\label{ch:la-hatred-and-disgust}
\hlchaptermodality{494}

\begin{chapteropening}
\textbf{By the end of this chapter:} \emph{I can talk about hatred, disgust and weariness.}

Complete the last lesson of chapter 494: I can talk about hatred, disgust and weariness.
\end{chapteropening}

\section[ōdī, ōdisse]{ōdī, ōdisse — to hate}

\label{lesson:LA-C494-odi}



\begin{quote}
Before the new one: say the Latin for unworthy, then the Latin for to be indignant.
\end{quote}

\subsection*{You'll want to know: ōdī, ōdisse}
\textbf{ōdī, ōdisse} — \textquotedblleft{}to hate\textquotedblright{}.

\begin{quote}
Fēlēs aquam ōdit.
\end{quote}

\begin{quote}
\emph{The cat hates water.}
\end{quote}

Perfect in form, present in meaning: \emph{ōdī} — I hate, \emph{ōderam} — I hated, \emph{ōderō} — I shall hate. Future participle \emph{ōsūrus}.

\begin{cousinweb}
\textbf{ōdī} is inherited, through Proto-Italic \emph{ōdai}, from a Proto-Indo-European root meaning \textquotedblleft{}to hate\textquotedblright{} (EtymDB). The noun \emph{odium} \textquotedblleft{}hatred\textquotedblright{}, source of English \emph{odium} and \emph{odious}, is built on it.
\end{cousinweb}

\begin{grammarlens}[title={What you've built}]
One more word for telling, explaining and describing at B1.
\end{grammarlens}

\subsection*{Your turn}
\begin{itemize}
\raggedright
  \item \emph{Say it:} \emph{ōdī, ōdisse}
  \item \emph{Say it:} \emph{ōdī, ōdisse}, once more
  \item \emph{Say it:} a sentence of your own with \emph{ōdī, ōdisse}
  \item \emph{Recall:} say the Latin for unworthy, then the Latin for to be indignant, then say \emph{ōdī, ōdisse} again
\end{itemize}

\begin{tcolorbox}[breakable,colback=teal!4,colframe=teal!35!black,title={Before you move on}]
What does ōdī, ōdisse mean? (\textquotedblleft{}To hate\textquotedblright{}.) Say it once more.
\end{tcolorbox}
\section[odiōsus, odiōsa, odiōsum]{odiōsus, odiōsa, odiōsum — hateful, tiresome}

\label{lesson:LA-C494-odiosus}



\begin{quote}
Before the new one: say the Latin for to be indignant, then the Latin for to hate.
\end{quote}

\subsection*{You'll want to know: odiōsus, odiōsa, odiōsum}
\textbf{odiōsus, odiōsa, odiōsum} — \textquotedblleft{}hateful, tiresome\textquotedblright{}.

\begin{quote}
Iter longum et odiōsum fuit.
\end{quote}

\begin{quote}
\emph{The journey was long and tiresome.}
\end{quote}

Built on \emph{odium} — hatred. In everyday Latin it often means simply annoying. English \emph{odious}.

\begin{cousinweb}
\textbf{odiōsus} is built on \emph{odium} \textquotedblleft{}hatred\textquotedblright{} with the ending \emph{-ōsus} \textquotedblleft{}full of\textquotedblright{} (Lewis \& Short; EtymDB). English \emph{odious} comes from it.
\end{cousinweb}

\begin{grammarlens}[title={What you've built}]
One more word for telling, explaining and describing at B1.
\end{grammarlens}

\subsection*{Your turn}
\begin{itemize}
\raggedright
  \item \emph{Say it:} \emph{odiōsus, odiōsa, odiōsum}
  \item \emph{Say it:} \emph{odiōsus, odiōsa, odiōsum}, once more
  \item \emph{Say it:} a sentence of your own with \emph{odiōsus, odiōsa, odiōsum}
  \item \emph{Recall:} say the Latin for to be indignant, then the Latin for to hate, then say \emph{odiōsus, odiōsa, odiōsum} again
\end{itemize}

\begin{tcolorbox}[breakable,colback=teal!4,colframe=teal!35!black,title={Before you move on}]
What does odiōsus, odiōsa, odiōsum mean? (\textquotedblleft{}Hateful, tiresome\textquotedblright{}.) Say it once more.
\end{tcolorbox}
\section[fastīdiō, fastīdīre]{fastīdiō, fastīdīre — to turn up one's nose at, feel disgust at}

\label{lesson:LA-C494-fastidio}



\begin{quote}
Before the new one: say the Latin for to hate, then the Latin for hateful.
\end{quote}

\subsection*{You'll want to know: fastīdiō, fastīdīre}
\textbf{fastīdiō, fastīdīre} — \textquotedblleft{}to turn up one's nose at, feel disgust at\textquotedblright{}.

\begin{quote}
Puer holera fastīdit.
\end{quote}

\begin{quote}
\emph{The boy turns up his nose at vegetables.}
\end{quote}

\emph{fastīdiō, fastīdīre, fastīdīvī, fastīdītum}, with the accusative.

\begin{cousinweb}
\textbf{fastīdiō} is built on \emph{fastīdium} \textquotedblleft{}loathing\textquotedblright{} (Lewis \& Short; EtymDB).
\end{cousinweb}

\begin{grammarlens}[title={What you've built}]
One more word for telling, explaining and describing at B1.
\end{grammarlens}

\subsection*{Your turn}
\begin{itemize}
\raggedright
  \item \emph{Say it:} \emph{fastīdiō, fastīdīre}
  \item \emph{Say it:} \emph{fastīdiō, fastīdīre}, once more
  \item \emph{Say it:} a sentence of your own with \emph{fastīdiō, fastīdīre}
  \item \emph{Recall:} say the Latin for to hate, then the Latin for hateful, then say \emph{fastīdiō, fastīdīre} again
\end{itemize}

\begin{tcolorbox}[breakable,colback=teal!4,colframe=teal!35!black,title={Before you move on}]
What does fastīdiō, fastīdīre mean? (\textquotedblleft{}To turn up one's nose at, feel disgust at\textquotedblright{}.) Say it once more.
\end{tcolorbox}
\section[fastīdium, fastīdiī]{fastīdium, fastīdiī — disgust, squeamishness; disdain}

\label{lesson:LA-C494-fastidium}



\begin{quote}
Before the new one: say the Latin for hateful, then the Latin for to turn up one's nose at.
\end{quote}

\subsection*{You'll want to know: fastīdium, fastīdiī}
\textbf{fastīdium, fastīdiī} — \textquotedblleft{}disgust, squeamishness; disdain\textquotedblright{}.

\begin{quote}
Aeger cibum cum fastīdiō spectat.
\end{quote}

\begin{quote}
\emph{The sick man looks at the food with disgust.}
\end{quote}

Neuter, second declension. Disgust of the stomach, and the scorn of a snob.

\begin{cousinweb}
\textbf{fastīdium} is a native Latin word; Lewis \& Short compares \emph{fastus} \textquotedblleft{}pride\textquotedblright{}, but the origin is not settled. English \emph{fastidious} comes from \emph{fastīdiōsus}, built on it.
\end{cousinweb}

\begin{grammarlens}[title={What you've built}]
One more word for telling, explaining and describing at B1.
\end{grammarlens}

\subsection*{Your turn}
\begin{itemize}
\raggedright
  \item \emph{Say it:} \emph{fastīdium, fastīdiī}
  \item \emph{Say it:} \emph{fastīdium, fastīdiī}, once more
  \item \emph{Say it:} a sentence of your own with \emph{fastīdium, fastīdiī}
  \item \emph{Recall:} say the Latin for hateful, then the Latin for to turn up one's nose at, then say \emph{fastīdium, fastīdiī} again
\end{itemize}

\begin{tcolorbox}[breakable,colback=teal!4,colframe=teal!35!black,title={Before you move on}]
What does fastīdium, fastīdiī mean? (\textquotedblleft{}Disgust, squeamishness; disdain\textquotedblright{}.) Say it once more.
\end{tcolorbox}
\section[taedet, taedēre]{taedet, taedēre — it wearies, it disgusts}

\label{lesson:LA-C494-taedet}



\begin{quote}
Before the new one: say the Latin for to turn up one's nose at, then the Latin for disgust.
\end{quote}

\subsection*{You'll want to know: taedet, taedēre}
\textbf{taedet, taedēre} — \textquotedblleft{}it wearies, it disgusts\textquotedblright{}.

\begin{quote}
Mē taedet longārum ōrātiōnum.
\end{quote}

\begin{quote}
\emph{I am weary of long speeches.}
\end{quote}

Impersonal: \emph{taedet, taedēre, taeduit} or \emph{taesum est}. The person in the accusative, the thing in the genitive: \emph{mē taedet labōris} — I am sick of the work. \emph{Taedium}, already known, is built on it.

\begin{cousinweb}
\textbf{taedet} is a native Latin impersonal verb, and the noun \emph{taedium} \textquotedblleft{}weariness\textquotedblright{} is built on it (Lewis \& Short); its deeper origin is not settled. English \emph{tedium} and \emph{tedious} come from \emph{taedium}.
\end{cousinweb}

\begin{grammarlens}[title={What you've built}]
One more word for telling, explaining and describing at B1.
\end{grammarlens}

\subsection*{Your turn}
\begin{itemize}
\raggedright
  \item \emph{Say it:} \emph{taedet, taedēre}
  \item \emph{Say it:} \emph{taedet, taedēre}, once more
  \item \emph{Say it:} a sentence of your own with \emph{taedet, taedēre}
  \item \emph{Recall:} say the Latin for to turn up one's nose at, then the Latin for disgust, then say \emph{taedet, taedēre} again
\end{itemize}

\begin{tcolorbox}[breakable,colback=teal!4,colframe=teal!35!black,title={Before you move on}]
What does taedet, taedēre mean? (\textquotedblleft{}It wearies, it disgusts\textquotedblright{}.) Say it once more.
\end{tcolorbox}
